\item The phrase ``larger steady-state error, particularly during the early
stages'' seems contradictory, since early-stage errors are typically transient
rather than steady-state. Please clarify whether ``transient error'' was intended
or if the statement refers to a different metric.

\textbf{Response:} The reviewer is correct: transient error was intended, and the
sentence has been corrected in the revised manuscript to read ``larger transient
error during the initial convergence phase''.

We take the opportunity to state why the effect is necessarily transient, which
the original phrasing obscured. The coupling term of the protocol is a signed
square root of the disagreement,
\[
  g_i \;=\; \alpha \sum_{j \in \mathcal{N}_i}
            -\operatorname{sign}(z_i - z_j)\sqrt{|z_i - z_j|},
\]
so a lost packet withholds a correction whose magnitude is
$\alpha\sqrt{|z_i - z_j|}$ and therefore scales with the prevailing disagreement.
With $\alpha = 0.02$, a single lost update costs $6.3\times10^{-2}$ in state units
when neighbours differ by $10$, but only $2.0\times10^{-3}$ when they differ by
$10^{-2}$. An identical packet-loss rate consequently produces a large error
contribution during the initial convergence phase and a vanishing one as the
agents approach agreement.

This is why the asymmetry between purely BLE links and bridge-to-BLE links is
visible early and not at steady state: the two link classes differ in loss rate,
but the cost of a loss decays with the disagreement it would have corrected. The
steady-state disagreement is instead set by the discretisation of the
finite-time term, which bounds the residual at $(\alpha d)^2$ for a node of
degree $d$, independent of the transport.
